\documentclass[11pt,a4paper]{article}
\usepackage[T1]{fontenc}
\usepackage[utf8]{inputenc}
\usepackage[french]{babel}
\usepackage{amsmath,amssymb}
\usepackage[margin=2.5cm]{geometry}
\usepackage[hidelinks,breaklinks=true]{hyperref}
\setlength{\emergencystretch}{2em}
\newcommand{\OO}{\mathcal O}
\newcommand{\PP}{\mathbf P}
\newcommand{\ZZ}{\mathbf Z}
\newcommand{\question}[1]{\begin{quote}\textbf{Question.}\enspace #1\end{quote}}
\newcommand{\reponse}{\noindent\textbf{Réponse.}\enspace}

\title{L'exemple d'ACL : cartes, obstructions et correction finie\\
\large Une discussion en questions et réponses}
\author{}
\date{27 octobre 2026}

\begin{document}
\maketitle

\noindent Ce document reprend les trois échanges consacrés à l'intuition
géométrique de la preuve cohomologique dans l'exemple d'ACL, au lemme de
correction finie et à leurs antécédents historiques. Les réponses sont
mises en forme pour la lecture, avec les références historiques réunies
à la fin. La distinction entre l'inspiration de la méthode et
l'attribution du lemme exact est conservée.

\section{L'intuition géométrique de l'exemple d'ACL}

\question{Pour ce que l'on a fait sur l'exemple de ACL, tu as une intuition
géométrique de la situation, en parlant de cartes, d'obstructions, etc.
Tu peux préciser avec une explication pour quelqu'un qui ne connaît pas
tout cela ? Je comprends l'algèbre, mais j'ai peu d'intuitions géométriques
et je voudrais en avoir plus.}

\reponse Oui. \textbf{L'image centrale est celle d'une famille de solutions
dont une partie reste près de l'origine, tandis qu'une autre peut partir
à l'infini.} La preuve construit des coordonnées qui permettent de
conserver la partie qui nous intéresse dans une algèbre finie.

Je vais expliquer les mots géométriques en les traduisant systématiquement
en algèbre.

Dans l'exemple d'ACL, on a
\[
\begin{aligned}
 p&=x-a+bxy+2bx^2,\\
 q&=y-b+ax^2+axy+by^2.
\end{aligned}
\]
On regarde $a,b$ comme des \textbf{paramètres}, et $x,y$ comme les inconnues.
Pour chaque valeur des paramètres, les équations donnent une configuration
de solutions. Cette configuration s'appelle la \emph{fibre}.
Algébriquement, prendre la fibre en $a=b=0$, c'est passer au quotient
par $(a,b)$.

Ici, cette fibre est simplement
\[
 x=0,\qquad y=0.
\]
Elle consiste donc en un seul point, l'origine. La difficulté est que
cette simplicité à l'origine des paramètres ne contrôle pas immédiatement
toutes les solutions pour les autres valeurs de $a,b$.

Pour voir ce phénomène dans un exemple plus petit, considérons
\[
 x+ax^2=0.
\]
Pour $a\neq0$, les deux solutions sont $x=0$ et $x=-1/a$.
Lorsque $a$ tend vers zéro, la seconde part à l'infini, tandis que la
première reste à l'origine. À $a=0$, l'équation affine ne montre plus
que $x=0$.

Dans cet exemple, $s=1+ax$ distingue exactement les deux branches :
\[
 s=1\text{ sur }x=0,\qquad s=0\text{ sur }x=-1/a.
\]
Inverser $s$ conserve la première et élimine la seconde. C'est une
première intuition du rôle de l'élément $s$ dans ZMT.

\subsection*{Rendre l'infini visible}

Pour rendre l'infini visible dans l'exemple d'ACL, on introduit les
coordonnées homogènes
\[
 [X:Y:Z],\qquad x=X/Z,\quad y=Y/Z.
\]
Deux triplets proportionnels représentent le même point. La région
$Z\neq0$ est notre plan affine habituel ; la région $Z=0$ est la droite
ajoutée à l'infini.

Les équations deviennent
\[
\begin{aligned}
 P&=XZ-aZ^2+bXY+2bX^2,\\
 Q&=YZ-bZ^2+aX^2+aXY+bY^2.
\end{aligned}
\]

\textbf{Au-dessus de $a=b=0$, elles deviennent $XZ=YZ=0$.}
Il y a donc deux morceaux disjoints :
\begin{itemize}
 \item si $Z\neq0$, on retrouve le point $[0:0:1]$ ;
 \item si $Z=0$, tous les points de la droite à l'infini satisfont les
 équations.
\end{itemize}
La fibre projective est ainsi le point qui nous intéresse \emph{et une
droite entière à l'infini}. Cette droite est un effet concret de
l'homogénéisation : lorsque les termes quadratiques disparaissent, les
équations deviennent toutes deux multiples de $Z$. Il ne faut pas
interpréter chacun de ses points comme une branche affine distincte.

Une \emph{carte} est simplement une région où l'on peut choisir des
coordonnées ordinaires :
\begin{center}
\renewcommand{\arraystretch}{1.3}
\begin{tabular}{lll}
\hline
Carte & Coordonnées & Ce qu'elle permet de voir\\
\hline
$Z\neq0$ & $X/Z,\ Y/Z$ & Le plan affine initial\\
$X\neq0$ & $Y/X,\ Z/X$ & L'infini où $X\neq0$\\
$Y\neq0$ & $X/Y,\ Z/Y$ & L'infini où $Y\neq0$\\
\hline
\end{tabular}
\end{center}
Algébriquement, une carte est obtenue en autorisant les dénominateurs
correspondants. L'intersection des cartes $X$ et $Z$ autorise les deux
dénominateurs. C'est pourquoi les polynômes de Laurent apparaissent
naturellement.

\subsection*{Ce qu'est une section}

\textbf{Une section globale est une formule donnée sur chaque carte,
avec accord sur les intersections.}

Il faut toutefois préciser quel genre de formule nous cherchons. Nous
voulons finalement trois fonctions régulières
\[
 s,\qquad sx,\qquad sy.
\]
Comme $x=X/Z$ et $y=Y/Z$, multiplier une fonction régulière $s$ par $x$
ou $y$ risque d'introduire des pôles à l'infini.

Pour contrôler les trois fonctions simultanément, nous cherchons une
expression $h$ telle que
\[
 Zh,\qquad Xh,\qquad Yh
\]
soient toutes régulières.

C'est exactement le rôle de $\OO(-1)$. Concrètement, une section de cet
objet s'écrit sur les trois cartes
\[
 \frac{u_X}{X},\qquad \frac{u_Y}{Y},\qquad \frac{u_Z}{Z},
\]
où chaque $u_L$ est une fonction régulière sur sa carte. Ces expressions
doivent être égales sur les intersections. Leur degré homogène est $-1$,
de sorte que les multiplier par $X,Y,Z$ donne des expressions de degré
zéro, donc de véritables fonctions.

Sur la fibre $a=b=0$, le choix est évident :
\[
 h_X=0,\qquad h_Y=0,\qquad h_Z=\frac1Z.
\]
Il vaut zéro sur le morceau à l'infini et $1/Z$ sur le point affine.
Ces prescriptions sont compatibles, puisque ces deux morceaux sont
disjoints.

Voilà la section que nous voulons relever lorsque $a,b$ ne sont plus
nuls.

\subsection*{Les désaccords entre cartes et les obstructions}

\textbf{L'obstruction mesure le défaut d'accord entre les formules
locales.}

Sur toute la famille, les mêmes trois expressions
\[
 z_*=(0,0,1/Z)
\]
ne se recollent plus nécessairement. Sur l'intersection des cartes $X,Z$,
par exemple, les deux prescriptions sont $0$ et $1/Z$.

L'opérateur de Čech $d$ enregistre simplement ces différences :
\[
 d(h_X,h_Y,h_Z)
   =(h_Y-h_X,\ h_Z-h_X,\ h_Z-h_Y).
\]
Ainsi
\[
 dz_*=(0,1/Z,1/Z).
\]

Un triplet est une section globale exactement lorsque ses différences
sont nulles. Et les différences vérifient toujours une relation :
\[
 (h_Z-h_Y)-(h_Z-h_X)+(h_Y-h_X)=0.
\]
C'est l'identité $d^2=0$, qui prend ici un sens très concret.

La \emph{cohomologie} intervient lorsqu'on demande si ces différences
peuvent être supprimées par des corrections locales soumises à certaines
contraintes.

Dans notre cas, on veut corriger sans changer la section modulo
$I=(a,b)$, ou même modulo $I^n$. Une erreur appartenant à $I^nK^1$ est
supprimable à cet ordre s'il existe une correction $v\in I^nK^0$ telle que
\[
 dz_*=dv.
\]
L'obstruction est la classe de $dz_*$ modulo ces corrections autorisées.

C'est une nuance importante : $dz_*$ est évidemment une différence de
cochaînes, mais \textbf{on demande une primitive dans $I^nK^0$}.
La primitive évidente $z_*$ n'y appartient pas.

\subsection*{Des approximations à tous les ordres}

Pourquoi disposons-nous d'approximations à tous les ordres ?

Dans l'algèbre affine $B$, la première équation donne
\[
 x=a-bxy-2bx^2\in IB.
\]
Sur l'intersection des cartes $X,Z$, $x$ est inversible. Donc
\[
 1\in I B[1/x].
\]
De même, la seconde équation donne $1\in I B[1/y]$.

Par conséquent, sur ces intersections, $1/Z$ appartient à chaque
puissance $I^n$. \textbf{Le désaccord disparaît modulo $I^n$, pour tout
$n$.}

Géométriquement : à chaque ordre infinitésimal autour de $a=b=0$, le
morceau affine reste séparé du morceau à l'infini.

Mais on ne peut pas conclure que le désaccord est nul. Dans $k[a,1/a]$,
par exemple, $1$ appartient à tous les idéaux $(a^n)$, sans être nul.
Cela explique pourquoi « on sait relever à tous les ordres » ne termine
pas la preuve.

Le lemme de correction finie est précisément le passage supplémentaire.
Le contrôle noethérien des obstructions fournit une relation finie
\[
 dz_*-\sum_i r_i\,dz_*=dv,
 \qquad r_i\in I,\quad v\in I^NK^0.
\]
On pose alors
\[
 h=z_*-\sum_i r_i z_*-v.
\]
Cette fois,
\[
 dh=0,\qquad h\equiv z_*\pmod I.
\]

Nous avons une section globale exacte, obtenue par une combinaison
finie. Aucune convergence de série n'est nécessaire.

\subsection*{Le retour à une algèbre finie}

On peut maintenant lire géométriquement la conclusion.

Les fonctions
\[
 Zh,\qquad Xh,\qquad Yh
\]
sont régulières sur toute la famille projective. Sur la carte affine,
elles s'écrivent
\[
 s,\qquad sx,\qquad sy.
\]
La réduction de $h$ donne $s\equiv1\pmod{(a,b)}$, donc l'ouvert $s\neq0$
contient la fibre affine qui nous intéresse. À l'infini, $Zh$ s'annule :
le facteur $Z$ élimine précisément cette région.

Il reste à obtenir les équations moniques. C'est le rôle du contrôle
fini des sections globales. Dans la preuve formalisée, on associe aux
sections des codes dans un module fini ; une dépendance noethérienne
entre les codes des puissances d'une section se traduit en une équation
monique pour cette section.

Ainsi $s,sx,sy$ sont entiers sur $A$. L'algèbre
\[
 D=A[s,sx,sy]
\]
est finie sur $A$, et après inversion de $s$,
\[
 x=\frac{sx}{s},\qquad y=\frac{sy}{s},
 \qquad D[1/s]=B[1/s].
\]

Le sens géométrique de ce résultat est précis : \textbf{la partie de
notre famille où $s\neq0$ se retrouve comme un ouvert d'une famille
finie sur les paramètres}. Le travail sur les cartes a construit des
coordonnées adaptées, $s,sx,sy$, dont le comportement reste contrôlé là
où les coordonnées initiales $x,y$ pouvaient partir à l'infini.

\section{Comment expliquer le lemme de correction finie ?}

\question{Et comment on explique le lemme de correction finie ?}

\reponse L'idée est la suivante : \textbf{on combine plusieurs
approximations pour que leurs défauts de recollement se compensent,
tout en conservant la section initiale modulo $I$}. La finitude garantit
qu'une telle compensation apparaît après un nombre fini d'étapes.

Écrivons cela avec seulement des modules et des applications linéaires :
\[
 K^0\xrightarrow{d}K^1\xrightarrow{d_1}K^2,
 \qquad d_1d=0.
\]
$K^0$ contient les collections de formules locales ; $d$ calcule leurs
désaccords. On cherche donc $h\in K^0$ avec $dh=0$.

Supposons que, pour chaque $n\geq1$, nous ayons une approximation $z_n$
telle que
\[
 dz_n\in I^nK^1,\qquad z_n\equiv z_1\pmod{IK^0}.
\]
Elles représentent toutes la même section initiale, et leurs défauts
sont d'ordres de plus en plus élevés.

\textbf{À l'ordre $n$, on oublie les erreurs que l'on sait déjà corriger.}
Plus précisément, on considère
\[
 M_n=
 \frac{\ker(d_1:I^nK^1\longrightarrow I^nK^2)}
      {d(I^nK^0)}.
\]
La classe
\[
 \beta_n=[dz_n]\in M_n
\]
est l'obstruction restante. Dire $\beta_n=0$, c'est dire qu'il existe
$v_n\in I^nK^0$ avec $dz_n=dv_n$ ; alors $z_n-v_n$ est la section
cherchée.

On ne sait pas nécessairement qu'une des $\beta_n$ s'annule.
\textbf{Mais il suffit qu'une obstruction soit combinaison des
précédentes, avec les bons coefficients.}

La relation utile est
\[
 \beta_N=\sum_{1\leq i<N}r_i\beta_i,
 \qquad r_i\in I^{N-i}.
\]
Par définition des classes, elle signifie qu'il existe $v\in I^NK^0$
tel que
\[
 dz_N-\sum_{1\leq i<N}r_i\,dz_i=dv.
\]
On pose simplement
\[
 \boxed{h=z_N-\sum_{1\leq i<N}r_i z_i-v.}
\]
Alors $dh=0$. Et puisque $N-i\geq1$, tous les $r_i$ appartiennent à $I$,
donc
\[
 h\equiv z_N\equiv z_1\pmod I.
\]
C'est tout le calcul du lemme.

Pourquoi la finitude donne-t-elle une relation avec \emph{ces puissances
précises de $I$} ?

On rassemble les obstructions dans un module gradué
\[
 M=\bigoplus_{n\geq0}M_n.
\]
Le degré garde la mémoire de l'ordre de l'erreur. Un coefficient dans
$I^j$ fait passer une obstruction d'ordre $i$ à une obstruction d'ordre
$i+j$. L'anneau qui organise cela est l'anneau de Rees
\[
 \mathcal R=\bigoplus_{j\geq0}I^jT^j.
\]
La lettre $T$ sert ici à enregistrer l'ordre.

Le contrôle noethérien approprié fournit une relation entre $\beta_N$
et les précédentes. En prenant sa composante de degré $N$, le
coefficient de $\beta_i$ appartient automatiquement à $I^{N-i}$.

\textbf{La graduation protège donc la section initiale.} Une dépendance
sans cette contrainte pourrait supprimer les erreurs tout en supprimant
aussi la section que nous voulions relever.

Dans ACL, c'est encore plus parlant : toutes les approximations utilisent
la même cochaîne $z_*=(0,0,1/Z)$. La construction devient
\[
 h=(1-r)z_*-v,\qquad r\in I,\quad v\in I^NK^0.
\]
On ajuste légèrement le coefficient de la prescription initiale, puis
on soustrait une correction locale. Les désaccords disparaissent
exactement, et la valeur modulo $I$ reste inchangée.

Le travail substantiel est donc de prouver ce \textbf{contrôle fini des
obstructions}. Une fois celui-ci disponible, le relèvement se fait par
cette combinaison linéaire finie, sans avoir à sommer une suite infinie
de corrections.

\section{Grothendieck et les antécédents historiques}

\question{C'est Grothendieck qui a trouvé cela ? Est-ce qu'il y avait
des précédents ?}

\reponse \textbf{Pour la stratégie cohomologique, oui ; pour notre lemme
exact de correction finie, je n'ai pas d'attribution historique
vérifiée.} Mon attribution précédente à Grothendieck était donc trop
globale.

Il y a plusieurs antécédents distincts.

\textbf{Zariski avait déjà relié les voisinages formels aux problèmes
de connexité.} Dans sa conférence de 1950, il explique que son théorème
de connexité repose sur une théorie de fonctions « holomorphes »
abstraites, développée dans son mémoire paru en 1951. Le lien entre
information infinitésimale et géométrie d'une fibre précède donc
Grothendieck~\cite[\S5]{Z1950}\cite{Z1951}.

\textbf{Le mécanisme algébrique de notre correction est très proche
d'Artin--Rees.} On rassemble les puissances d'un idéal dans un anneau
gradué, puis on utilise la noethérianité pour contrôler tous les ordres
à partir de données finies. Cette méthode appartient aux travaux de
Rees des années 1950 ; son article \emph{Two classical theorems of ideal
theory} paraît en 1956~\cite{Rees}. Notre usage de la graduation pour
obtenir les coefficients $r_i\in I^{N-i}$ relève de cette même idée.

\textbf{Grothendieck en fait une théorie cohomologique générale.}
Le théorème des fonctions formelles compare la cohomologie avant et
après complétion ; dans EGA III, il conduit au théorème de connexion
puis à une forme de ZMT. C'est cette organisation de la preuve qui a
guidé notre démarche~\cite[\S\S4.1--4.4]{EGAIII}.

Et \textbf{Serre intervient dans la preuve elle-même} : Illusie précise
que la preuve originale de Grothendieck n'a pas été publiée, et que
celle d'EGA III, 4.1.7--4.1.8, suit un argument dû à
Serre~\cite[remarque 2.3(c)]{Illusie}.

Concernant précisément
\[
 h=z_N-\sum_{i<N}r_i z_i-v,
\]
je ne prétendrais donc ni que Grothendieck a énoncé ce lemme sous cette
forme, ni que cette formulation est nouvelle. \textbf{Nous avons isolé
un mécanisme fini au sein de cette tradition}, en demandant seulement
de relever la section modulo $I$, et en utilisant une relation entre
les obstructions rencontrées plutôt qu'un théorème complet de
comparaison avec les complétés. C'est cette réduction des besoins qui
rend l'argument particulièrement adapté à notre formalisation
constructive.

\begin{thebibliography}{9}

\bibitem{Z1950}
O. Zariski,
conférence au Congrès international des mathématiciens de 1950,
\emph{Proceedings of the International Congress of Mathematicians,
Cambridge, Massachusetts, 1950}, vol.~II,
\S5, « Holomorphic functions and the principle of degeneration », p.~84.
\href{https://www.mathunion.org/fileadmin/ICM/Proceedings/ICM1950.2/ICM1950.2.ocr.pdf}
{Actes du congrès, volume II (PDF)}.

\bibitem{Z1951}
O. Zariski,
\emph{Theory and Applications of Holomorphic Functions on Algebraic
Varieties over Arbitrary Ground Fields},
Memoirs of the American Mathematical Society, no.~5, 1951.
\href{https://doi.org/10.1090/memo/0005}{DOI : 10.1090/memo/0005}.

\bibitem{Rees}
D. Rees,
\emph{Two classical theorems of ideal theory},
Proceedings of the Cambridge Philosophical Society, 52 (1956),
p.~155--157.
\href{https://doi.org/10.1017/S0305004100031091}
{DOI : 10.1017/S0305004100031091}.

\bibitem{EGAIII}
A. Grothendieck, avec la collaboration de J. Dieudonné,
\emph{Éléments de géométrie algébrique III : Étude cohomologique des
faisceaux cohérents, première partie},
Publications mathématiques de l'IHÉS, 11 (1961),
\S\S4.1--4.4.
\href{https://wstein.org/sga/circle/EGAIII.1.pdf}{Texte intégral (PDF)}.

\bibitem{Illusie}
L. Illusie,
\emph{Grothendieck's existence theorem in formal geometry},
notes d'août 2003, remarque 2.3(c).
\href{https://indico.ictp.it/event/a0255/session/28/contribution/17/material/0/0.pdf}
{Notes disponibles à l'ICTP (PDF)}.

\end{thebibliography}

\end{document}
